% Shared preamble + layout macros for Ratio competitive component specs.
% Each <slug>.tex \input{}s this, then opens \begin{document}. Compiled
% per-component to its own PDF by rules_tectonic (see spec.bzl / BUILD).
\documentclass[11pt,letterpaper]{article}
\usepackage[T1]{fontenc}
\usepackage[utf8]{inputenc}
\usepackage{lmodern}
\usepackage{microtype}
\usepackage[margin=0.9in]{geometry}
\usepackage{xcolor}
\usepackage{graphicx}
\usepackage{booktabs}
\usepackage{array}
\usepackage{colortbl}
\usepackage{pifont}
\usepackage{enumitem}
\usepackage[most]{tcolorbox}
\usepackage{hyperref}

% ── Brand palette (shared with the Ratio papers) ────────────────────
\definecolor{ratioInk}{HTML}{0F2418}
\definecolor{ratioAccent}{HTML}{1B6440}
\definecolor{ratioMid}{HTML}{679177}
\definecolor{ratioRule}{HTML}{C3D6BF}
\definecolor{ratioTint}{HTML}{E6EFE2}
\definecolor{posGreen}{HTML}{1F6B45}
\definecolor{negGray}{HTML}{666666}

\hypersetup{colorlinks=true, linkcolor=ratioAccent, citecolor=ratioAccent, urlcolor=ratioAccent}
\setlist[itemize]{leftmargin=1.3em, itemsep=2.5pt, topsep=2pt}
\pagestyle{plain}
\setlength{\parindent}{0pt}
\setlength{\parskip}{4pt}

\newcommand{\yes}{{\color{posGreen}\ding{51}}}
\newcommand{\no}{{\color{negGray}\ding{55}}}

% ── Header banner: \spechead{Title}{slug} ───────────────────────────
\newcommand{\spechead}[2]{%
  \noindent\begin{tcolorbox}[colback=ratioInk, colframe=ratioInk, arc=4pt,
      left=14pt, right=14pt, top=11pt, bottom=11pt, boxrule=0pt]
    {\sffamily\footnotesize\color{ratioRule}
      RATIO COMPONENT SPEC \;\textbullet\; benchmarked to the wealth-tech stack
      \;\textbullet\; \texttt{#2}}\\[3pt]
    {\sffamily\bfseries\LARGE\color{white} #1}
  \end{tcolorbox}\par\vspace{8pt}}

% ── Section heading: \specsection{Name} ─────────────────────────────
\newcommand{\specsection}[1]{%
  \vspace{6pt}{\sffamily\large\bfseries\color{ratioAccent} #1}\par
  \vspace{1pt}{\color{ratioRule}\rule{\linewidth}{1pt}}\par\vspace{3pt}}

% ── Highlighted "Ratio mapping" box ─────────────────────────────────
\newtcolorbox{ratiobox}{colback=ratioTint, colframe=ratioMid, boxrule=0.8pt,
  arc=3pt, left=9pt, right=9pt, top=7pt, bottom=7pt}

% ── Status footer box ───────────────────────────────────────────────
\newtcolorbox{statusbox}{colback=ratioRule!35, colframe=ratioRule, boxrule=0.6pt,
  arc=2pt, left=8pt, right=8pt, top=5pt, bottom=5pt}

% ── "Advertised by" two-column table helper ─────────────────────────
\newcommand{\advrow}[2]{\textbf{#1} & #2 \\}

\begin{document}
\spechead{Risk Analytics \& Scenario Modeling}{risk-analytics}

Forward-looking analysis: risk metrics, stress tests, and what-if scenario
projection over current portfolios.

\specsection{Advertised by (benchmark sources)}
\begin{tabular}{@{}p{0.34\linewidth}p{0.60\linewidth}@{}}
\advrow{Orion}{Risk Intelligence --- ``understand risk, uncover opportunity''}
\advrow{Addepar}{Navigator --- advanced scenario modeling / project future outcomes}
\advrow{SS\&C Black Diamond}{Monte Carlo simulations (investment management)}
\advrow{Envestnet Tamarac}{Risk/analytics within Reporting}
\end{tabular}

\specsection{Advertised capabilities (feature level)}
\begin{itemize}
  \item Portfolio risk metrics and exposure analysis.
  \item Scenario / what-if modeling to project future outcomes.
  \item Monte Carlo simulation for goal and outcome probabilities.
  \item Stress testing; opportunity identification across portfolios.
\end{itemize}

\begin{ratiobox}
\textbf{Ratio mapping.} Scenario modeling runs candidate states through a
sandboxed instance of the proven kernel, seeded with typed market facts ---
so projections inherit the engine's exactness and the inputs' provenance.
Risk computation is a read-only analytical layer (read-only Rust extensions) over
ledger projections; it never mutates the book. ``Which price set drove this
projection'' is answerable by hash.
\end{ratiobox}

\specsection{Conservation \& trust placement}
\textbf{Analytical tier + simulation over the trusted core.} Read-only;
provenance-bound to fact sets.

\specsection{References}
\footnotesize
\begin{itemize}
  \item \url{https://orion.com/advisor-tech}
  \item \url{https://addepar.com/wealth-management}
\end{itemize}

\begin{statusbox}
\textbf{Status:} scaffold --- generated from a June 2026 competitive scrape.
Expand into a full Ratio component spec. Capabilities reflect competitors' public
marketing, not Ratio commitments.
\end{statusbox}
\end{document}
